%%%%%%%%%%%%%%%%%%%%%%%%%%%%%%%%%%%%%%%%%%%%%%%%%%%%%%%
%                                                     %
% Cette version repart de 03_Draft\intro-short-V1.tex %
%                                                     %
%%%%%%%%%%%%%%%%%%%%%%%%%%%%%%%%%%%%%%%%%%%%%%%%%%%%%%%

% NB: Je supprime toutes les (grosses) parties qui étaient en commentaire dans la version précédente pour qu'on y voit plus clair.

Since \cite{Jacobson1993} and \cite{Farber1993}, the literature on job displacement has documented long-run earnings losses following layoffs. This empirical regularity highlights the costs associated with the reallocation of factors of production. While this reallocation process is at the core of productivity growth in modern economies \citep{foster2001aggregate,hsieh2009misallocation}, it can entail substantial {private} costs.

Recent contributions on job displacement emphasize the role of firms in explaining the private cost of this reallocation. In particular, these studies show that displaced workers reallocate to firms with lower wage premia, measured \textit{à la} \cite{Abowd1999} (henceforth AKM). Yet, little is known regarding the productivity and the labor share of those (destination) firms. Productivity and wage premia being positively related across firms \citep[e.g.][]{card2016bargaining}, the reallocation of displaced workers to lower-wage firms appears likely to reflect movements toward lower productivity firms. But it could instead be driven by reemployment at low labor share, high productivity firms \citep{autor2020fall,aghion2022theory}. Both possibilities are compatible with the documented earnings losses but have different implications regarding \textit{i)} the ``productivity-enhancing'' nature of labor reallocation in the wake of downsizing events \citep{foster2016reallocation}, and \textit{ii)} whether earnings losses reflect a decline in workers' negotiation power or reallocation to firms generating less surplus per worker.%\footnote{Here, we simply define loss in negociation power as being re-employed by firms with unfavorable wage policy \textit{given}, as opposed to \textit{because} of, their productivity.}

In this paper, we shed new light on these questions by studying jointly the wage policy and productive performance of the firms in which displaced workers are re-employed. We proceed in three steps. First, using administrative data from France, we study a large set of workers laid-off for economic reasons, and find evidence of large and persistent earnings losses---about a third of pre-displacement earnings 6 years post-displacement (\textcolor{red}{could add results}) \textcolor{orange}{[32\%: $-7{,}961$ euros at $t^D+6$ (Appendix Table A7, col. 2) over the 24,848-euro baseline of Table 1; $-0.322$ in logs]}. Using detailed administrative data on hours worked, we find that working time explains almost all of the short run losses (1 to 2 years post-displacement), but tends to recover in the medium run (5 years post-displacement). By contrast, losses in hourly wage rate, while initially smaller, show no sign of recovery. These persistent losses in hourly wages are explained to a substantial extent (43\% -> \textcolor{red}{37.1\% 5 years after displacement}) \textcolor{orange}{[37\% OK: Table 2, col. 7 at $t^D+5$ = 0.371]} by declines in firm wage premia. Our findings add to an active literature assessing the importance of lost wage premia in shaping the cost of displacement \citep{Gulyas2019,Lachowska2020,moore2019firm,Schmieder2020,fackler2017explaining,bertheau_etal_2023,helm2022}.\footnote{The contribution of firm-specific wage premia to earnings losses remains debated, with some studies finding large effects (e.g. \citealp{Gulyas2019,Schmieder2020}) and others more modest ones (e.g. \citealp{Lachowska2020}); see Section \ref{sec:FEtheory} for further discussion.}

% Alternative with more references, including Bertheaud qu'on nous reprochait de ne pas assez citer, mais du coup comme on va le remettre dans le main, on peut sans doute s'en passer dans la footnote: 
%\footnote{The contribution of firm-specific wage premia to earnings losses remains debated, with some studies finding large effects (e.g. \citealp{Gulyas2019,Schmieder2020}) and others more modest ones (e.g. \citealp{Lachowska2020}); cross-country evidence is provided by \citealp{bertheau_etal_2023}. See Section \ref{sec:FEtheory} for further discussion.}

% original footnote:
%\footnote{The role of firm-level wage premium in overall earnings and wage rate losses of displaced workers is an active topic of research. \cite{Gulyas2019} and \cite{Schmieder2020} find firm premium important for shaping earnings losses, consistent with the work of \cite{goldschmidt2017rise} on outsourced workers. Conversely, \cite{Lachowska2020} find this loss to be of modest importance. \cite{bertheau_etal_2023} provide cross-country evidence. Further literature discussion is in section \ref{sec:FEtheory}.}

In a second step, we investigate the patterns of workers' reallocation in terms of firms' characteristics, such as productivity or labor share. We show that the reallocation of workers to lower-wage-premium firms does not reflect a reallocation towards low-productivity firms. Instead, despite the positive correlation between productivity and wage premia in the cross-section of firms, we find that destination firms tend to be more productive but pay lower premia.\footnote{This result is compatible with \cite{bell_bukowski_machin_2024} who show that the pass-through from rents to wages tends to decline over time, using UK data.} Those firms consequently have much lower labor share. Crucially, this analysis does not assume that displaced and control workers initially work in similar firms. Rather, it exploits exogenous displacement from the workers' perspective, and studies its consequences in terms of reallocation. Furthermore, we use pre-displacement measures of firms' characteristics to ensure they do not reflect post-layoff dynamics. Overall, we view this pattern as suggesting that losses in premia reflect a decline in bargaining power as opposed to difficulties to access jobs in productive firms. Finally, we strengthen the case for a loss in bargaining power by examining two additional qualitative outcomes that capture the state of labor relations at the firm level. We show that displaced workers are reemployed by firms that \textit{i)} are less likely to reach firm-level collective wage agreements and \textit{ii)} have lower turnout in professional elections. 

We contribute to the literature in three main ways. First, thanks to an administrative definition of economic layoffs, we are able to focus on a larger set of involuntary separations and a more representative set of firing firms (``origin firms'' henceforth) than what is captured in the extant literature relying on mass layoffs. Crucially, we show that our set of origin firms are more representative in terms of wage premia, so that losses in premia following displacement are less likely to be driven by mean-reversion.\footnote{See \cite{Schmieder2020} section IV.A on mean reversion in establishment premia loss. } While the qualitative reallocation pattern documented above---displaced workers moving to firms with higher productivity but lower wage premia and lower labor shares---is preserved when we implement a mass-layoff design, we find that the share of the hourly wage losses explained by wage premium losses is larger in the mass-layoff sample, presumably because that design oversamples higher-premium origin firms. Second, %taking advantage of the possibility offered by French data to 
linking matched employer-employee data with exhaustive accounting data on firm performance, we compare losses in wage premia with measures of productivity, surplus per worker and labor share of the destination firms. To the best of our knowledge, we are the first to carry out this data intensive exercise when studying job displacement.\footnote{\cite{haltiwanger2018moves} study workers flows and firm productivity but focus on the reallocation of workers through job-to-job moves and do not study AKM-type wage premia. Closer to our exercise are \citet{haltiwanger2018cyclical} and \citet{bertheau_vejlin_2023}, who study worker flows across firms classified by average wage and labor productivity over the business cycle.} It reveals that destination firms are both more productive and lower labor share firms.
This finding speaks to random-search models of the labor market \citep{Lentz2024}. In these models, workers climb a job ladder via on-the-job search, moving to higher-surplus firms or renegotiating their wage upward when outside offers arrive, while displacement pushes them back down the ladder. The AKM firm wage premium $\psi_f$ identifies a firm-specific component of wages common to all workers of the firm. Structural displacement models in this class such as \citet{krolikowski2017job} and \citet{Audoly_2024} attribute the firm-side wage losses to reallocation toward less productive firms. Our findings point to a more nuanced picture: re-employed workers move \emph{up} the productivity distribution but \emph{down} the wage-premium distribution. The estimation and calibration of such models would therefore benefit from using direct productivity data and contrasting their results against the empirical pattern of reallocation along the productivity distribution. Finally, we assemble new data to corroborate the idea that earnings losses associated with displacement reflect a loss in collective bargaining. We build original firm-level indicators capturing the state of labor relations in each employer, in the pre-displacement period. The use of firm-level data (as opposed to sector-level proxies for collective bargaining) is crucial as we show that our effects are driven by within-sector dynamics. We provide evidence that displaced workers are reallocated to employers that are less likely to successfully conclude firm-level collective wage agreements and where workers have a lower participation rate during professional elections---arguably reflecting a lower degree of organization of the workforce.


%\todo[inline]{Emphasize novelty: lower issue of mean reversion because more representative of broad set of separations.}

